\documentclass[a4paper,12pt,fleqn]{article}
\usepackage[unicode,colorlinks=false]{hyperref}
\def\UrlBreaks{\do\/\do-\do\&\do.\do:}

\usepackage[utf8x]{inputenc}
%

\usepackage[L7x]{fontenc}
\usepackage{times}
\usepackage{ucs}
%\usepackage{microtype}
%\DisableLigatures{encoding = *, family = *}
 %package to switch the language
\usepackage{etoolbox}

  %set up of the page margins
\usepackage[top=2cm, bottom=2cm, left=3cm, right=1.5cm]{geometry}

 %1.1 line spacing
\linespread{1.1}


  %page numbering at the right side
\usepackage{fancyhdr}
\pagestyle{fancyplain}
\fancyhf{}
\renewcommand{\headrulewidth}{0pt} 
\fancyhfoffset[RO]{0cm}

  %to number at the bottom (exchange lines to number at the top)
\rfoot{\thepage}
  %\rhead{\thepage} %

% \usepackage[usenames,dvipsnames]{pstricks}
\urlstyle{same}
\hypersetup{
%  citecolor=Blue,
%  linkcolor=Blue,
%  urlcolor=Blue
pdfborder={0 0 0 }
}

 %for includegraphics
\usepackage{graphicx}

\usepackage[toc,page]{appendix}


\usepackage{caption}

 %for source codes
\usepackage{listings}
\lstset{commentstyle=\color{red},xleftmargin=10pt, framexleftmargin=6pt, numbersep=1mm, frame=single, numbers=left,numberstyle=\footnotesize,extendedchars=\true, inputencoding=utf8x,basicstyle=\footnotesize,extendedchars=true,
 keywordstyle=\color{black}\bfseries, breaklines=true, breakautoindent=true,framesep=8pt,linewidth=0.95\textwidth
}

 %for algorithms
\usepackage{algorithm}
\usepackage{algorithmic}
 %instead of the above two packages we can use algorithms2e
 %\usepackage[boxed,linesnumbered,vlined,slide]{algorithm2e}

 %special symbols
\usepackage{amsfonts}
%\usepackage{amssymb} %redundant
\usepackage{amsmath}

 %for theorem like environments
\usepackage{amsthm}

 \usepackage{datetime}
 \renewcommand{\dateseparator}{--}


% SI system units
\usepackage{siunitx}
\sisetup{detect-all}
% Problem with fonts \SI{x.xx}{\micro\metre}, solved with updmap-sys --enable Map=utm.map
\renewcommand{\sfdefault}{uhv}
%\renewcommand{\rmdefault}{utm}  %commented due to text missformating
\usepackage{stix}
\renewcommand{\ttdefault}{ucr}

% List management (itemize, etc.)
\usepackage{enumitem}

\newcommand*{\urlw}[1]{\href{#1}%
            {\nolinkurl{#1}}}

\numberwithin{equation}{section}


%%%%%%%%%%% lino įdėta
%
\usepackage{pifont,mdframed}

\newenvironment{warning}
  {\par\begin{mdframed}[linewidth=2pt,linecolor=red]%
    \begin{list}{}{\leftmargin=1cm
                   \labelwidth=\leftmargin}\item[\Large\ding{43}]}
  {\end{list}\end{mdframed}\par}


\newtoggle{inLithuanian}
\settoggle{inLithuanian}{false}

\newtoggle{needPreface}
\settoggle{needPreface}{false}

\newtoggle{signaturesOnTitlePage}
\settoggle{signaturesOnTitlePage}{false}

\lstset{language=PHP}


\theoremstyle{definition}
\newtheorem{definition}{\keyWordDefinition}
\newtheorem{example}{\keyWordExample}
\def\QED{\unskip\nobreak\hfill\kern5pt$\Box$}

\iftoggle{inLithuanian}{
%\usepackage[L7x]{fontenc}
\usepackage[english,lithuanian]{babel}

\newcommand{\todayiso}{\the\year \dateseparator \twodigit\month \dateseparator \twodigit\day}


\renewcommand{\today}{\number\year\space m. \space \ifcase\month\or
  sausio\or vasario\or kovo\or balandžio\or gegužės\or birželio\or
  liepos\or rugpjūčio\or rugsėjo\or spalio\or lapkričio\or
  gruodžio\fi
  \space\number\day\space d.}


 \usepackage{tocloft}
 \renewcommand\cftsecaftersnum{.} 
 \renewcommand\cftsubsecaftersnum{.} 
 \renewcommand\cftsubsubsecaftersnum{.}

 \usepackage{VUMIFKK}

 \DeclareCaptionLabelFormat{captionlt}{#2 #1}
   %smth is not fine with algorithms 
 \DeclareCaptionLabelFormat{captionltalg}{#2 #1 algoritmas}

 \usepackage{indentfirst}
 \renewcommand{\appendixtocname}{Priedai}
 \renewcommand{\appendixpagename}{Priedai}
 \renewcommand{\contentsname}{Turinys} 

 \renewcommand{\lstlistingname}{išeities kodas}
 \renewcommand{\figurename}{pav}
 \renewcommand{\tablename}{lentelė}


 \captionsetup*[lstlisting]{   
 labelsep=period,labelformat=captionlt
 }
 \captionsetup*[figure]{   
% labelsep=period,
 labelsep=space, %babel redefines pav to pav.
 labelformat=captionlt
 }
 \captionsetup*[table]{   
  labelsep=period,
  labelformat=captionlt
 }
 \renewcommand{\algorithmicrequire}{\textbf{Įvestis:}}
 \renewcommand{\algorithmicensure}{\textbf{Išvestis:}}

 \captionsetup*[algorithm]{   
 labelsep=period,labelformat=captionltalg
 }

\renewcommand{\thmhead}[3]{#2 #1#3}

}
{
%\usepackage[OT1,T1]{fontenc}
%\usepackage[L7x]{fontenc}



\usepackage[english]{babel}
\newcommand{\todayiso}{\twodigit\month \dateseparator \twodigit\day \dateseparator \the\year}
 \captionsetup*[algorithm]{   
 labelsep=period
 }
\captionsetup*[lstlisting]{   
 labelsep=period
 }
 \captionsetup*[figure]{   
 labelsep=period
 }
 \captionsetup*[table]{   
 labelsep=period
 }


}

%some kywords
 \def\keywordAbstract{\iftoggle{inLithuanian}{Santrauka}{Abstract}}
 \def\keywordAbstractOther{\iftoggle{inLithuanian}{Summary}{Santrauka}}
 \def\keyWordIntroduction{\iftoggle{inLithuanian}{Įvadas}{Introduction}}
 \def\keyWordConclusions{\iftoggle{inLithuanian}{Išvados ir rekomendacijos}{Conclusions and Recommendations}}
 \def\keyWordAIdeclaration{\iftoggle{inLithuanian}{Dirbtinio intelekto naudojimo deklaracija}{Declaration of Artificial  Intelligence}}

 \def\keyWordPreface{\iftoggle{inLithuanian}{Pratarmė}{Preface}}
 \def\keyWordAppendice{\iftoggle{inLithuanian}{Priedas}{Appendix}}
 \def\keyWordSignature{\iftoggle{inLithuanian}{parašas}{signature}}
 \def\keyWordDefinition{\iftoggle{inLithuanian}{apibrėžimas}{Definition}}
 \def\keyWordExample{\iftoggle{inLithuanian}{pavyzdys}{Example}}

\newcommand{\bothabstracts}[3]{
\setcounter{secnumdepth}{0}
\newpage
\hspace{2cm}
{\centering{\section{\keywordAbstract}}}

#1
\newpage
\hspace{2cm}
{\centering \section{\keywordAbstractOther}}

\begin{center}{\textbf{#2} }\end{center}

 #3
\setcounter{secnumdepth}{3}
}

 %non-numbered sections: #1 param: for labeling sec:#1, #2 -section title
\newcommand{\sectionWithoutNumber}[2]{\newpage
%\hspace{2cm}
\section*{#1}
\label{sec:#2}
\addcontentsline{toc}{section}{\nameref{sec:#2}}%{#3}
 }



\newcommand{\referenceSources}[1]{
\newpage
\cleardoublepage
\phantomsection
\iftoggle{inLithuanian}{
 \renewcommand{\refname}{Literatūros šaltiniai}

 \addcontentsline{toc}{section}{Literatūros šaltiniai}
 \markboth{\refname}{Literatūros šaltiniai}
 }
{

\addcontentsline{toc}{section}{References}
\markboth{References}{References}
}

\bibliographystyle{plain}
\bibliography{#1}
}



 \newcommand\authorsignature[1]{
\begin{flushright}
 \begin{minipage}[b]{0.45\textwidth}
  \centering
  \rule{\textwidth}{0.5pt}\\
   #1
  \end{minipage}
\end{flushright}
 }




 \newcommand\authorsignatures[5]{%
   \vspace{1cm}
   \authorsignature{#1}
   \ifstrequal{#2}{}{}{\vspace{0.3cm}
     \authorsignature{#2}
     \ifstrequal{#3}{}{}{\vspace{0.3cm}
      \authorsignature{#3}
      \ifstrequal{#4}{}{}{\vspace{0.3cm}
        \authorsignature{#4}
        \ifstrequal{#5}{}{}{\vspace{0.3cm}
         \authorsignature{#5}       
        }
      }
    }
} 
}

\newcommand{\authortitle}{
\iftoggle{signaturesOnTitlePage}{
\tiny{\keyWordSignature}
}{}
}

\newcommand{\depttitlepage}[8]
{
\thispagestyle{empty}
\begin{center}

\includegraphics[width=2cm]{jb_VU_zenklas}

%\vspace{-1cm}

\iftoggle{inLithuanian}
{ 
  VILNIAUS UNIVERSITETAS\\
  MATEMATIKOS IR INFORMATIKOS FAKULTETAS\\
  INFORMATIKOS INSTITUTAS\\
  <<KOMPIUTERINIO IR DUOMENŲ MODELIAVIMO KATEDRA>>
}
{
  VILNIUS UNIVERSITY \\
  FACULTY OF MATHEMATICS AND INFORMATICS \\
  INSTITUTE OF COMPUTER SCIENCE\\
  <<DEPARTMENT OF COMPUTATIONAL AND DATA MODELING>>
}

\vspace{5cm}

#1\\
\vspace{0.5cm}
\textbf{\Large #2}
\end{center}

\vspace{5cm}


\begin{center}
\begin{minipage}{0.5\textwidth}
 \centering
\iftoggle{inLithuanian}
{
 \ifstrequal{#3}{}{}{Atliko:\\[5pt]}
}
{
\ifstrequal{#3}{}{}{Done by:\\[5pt]}
}
%\noindent
\begin{tabular}{@{}c@{}}%\setlength\tabcolsep{0pt}
\ifstrequal{#3}{}{}{#3\authortitle\\[5pt]}
\ifstrequal{#4}{}{}{#4\authortitle\\[5pt]}
\ifstrequal{#5}{}{}{#5\authortitle\\[5pt]}
\ifstrequal{#6}{}{}{#6\authortitle\\[5pt]}
\ifstrequal{#7}{}{}{#7\authortitle\\}
\end{tabular}

\end{minipage}
\end{center}

\vspace{0.5cm}
\begin{center}
\begin{minipage}{0.5\textwidth}
 \centering

\ifstrequal{#8}{}{}
{

\iftoggle{inLithuanian}
{
Vadovas:
}
{
Supervisor:
}

#8

}


\end{minipage}
\end{center}


\vfill

\begin{center}
Vilnius\\
\the\year
\end{center}

\iftoggle{needPreface}{
 \sectionWithoutNumber{\keyWordPreface}{preface}
\input{preface.tex}

\iftoggle{inLithuanian}
{
\vspace{\baselineskip}\hfill
\today
}
{
 \vspace{\baselineskip}\hfill \today
}

 \vspace{5cm}

\iftoggle{signaturesOnTitlePage}{}
{
\authorsignatures{#3}{#4}{#5}{#6}{#7}
}
}{}
\newpage
}


\begin{document}
 % #1 -report type, #2 - title, #3-7 students, #8 - supervisor
 \depttitlepage{Semester Project}{Fixing the Mathematical Foundations of the Lietuva Travel Category-Theoretic Engine}{Adrián Guinea Pérez}{}{}{}{}{Prof. I. Grinis}

\tableofcontents

\sectionWithoutNumber{Glossary of Terms}{keywords}{}
\begin{center}
\small
\begin{tabular}{|l|p{10cm}|}\hline
\textbf{Term} & \textbf{Meaning in this report} \\\hline
Category & Objects + composable, associative morphisms with identities (Sec.~\ref{sec:background}). \\\hline
Functor & Structure-preserving map between categories; \texttt{RouteList::fmap}. \\\hline
Natural transformation & A function-indexed family of morphisms satisfying a naturality square; \texttt{RouteList::headOption}. \\\hline
Monad & Functor with \texttt{pure}/\texttt{bind} satisfying three laws; \texttt{Option}, \texttt{Either}. \\\hline
Monoid & Associative operation with identity; \texttt{TripMetrics} under \texttt{append}. \\\hline
Free category $F(G)$ & Category generated by a directed graph $G$; objects = vertices, morphisms = paths. \\\hline
Kleisli composition & Chaining functions of type $a \to M(b)$ through a monad's \texttt{bind}. \\\hline
Short-circuit & An \texttt{Either::bind} on \texttt{Left} propagating \texttt{Left} unchanged. \\\hline
\end{tabular}
\end{center}

 %abstract (English only)
\sectionWithoutNumber{\keywordAbstract}{abstract}
This report documents the correction of three mathematical defects identified by Prof. I.~Grinis in the \emph{Lietuva Travel} categorical travel-planning engine (PHP~8.1+). Although the project's self-reported test suite claimed all its category-theoretic laws verified, three real violations survived: (1) the \texttt{Option} monad's \texttt{pure} constructor conflated the falsy value \texttt{false} with absence of a value, breaking the monad's right-identity law for \texttt{Some(false)}; (2) the \texttt{TripMetrics} monoid rounded its cost to two decimals after every \texttt{append()}, which made the monoid's associativity law fail for sub-cent amounts; and (3) \texttt{Itinerary::advance()} silently discarded a failed path composition in the induced free category $F(G)$, letting an invalid itinerary succeed instead of failing. Each defect is analysed at its root cause and fixed with the smallest change that restores the law. The project's additional invariant suite, \texttt{LietuvaTravel\_regression\_checks.php}, independently exercises the \texttt{Option} monad, the \texttt{TripMetrics} monoid, and the free-category path/itinerary machinery under edge-case inputs (nulls, sparse arrays, mismatched graph edges), reporting 10/10 checks passing.


 %Introduction section: label is sec:intro
\sectionWithoutNumber{\keyWordIntroduction}{intro}
\emph{Lietuva Travel} is a PHP~8.1+ travel-planning application whose core logic (\texttt{src/CategoryTheory.php}) is explicitly built on category-theoretic structures, following the constructions of Barr and Wells~\cite{BarrWells1990}: a \texttt{Functor}/\texttt{Monad} interface pair in the style of Moggi's and Wadler's monadic treatment of computation~\cite{Moggi1991,Wadler1992}, the \texttt{Option} and \texttt{Either} monads, the \texttt{RouteList} functor with a natural transformation into \texttt{Option}, the \texttt{TripMetrics} monoid for associative cost aggregation, and a directed graph whose paths are lifted into a free category $F(G)$ (Prof.~I.~Grinis' model), where edge composition is interpreted as Kleisli composition in the \texttt{Either} monad.

The project ships its own formal verification suite (\texttt{tests/category\_laws\_test.php}), which prints a ``15/15 laws verified'' badge. Prof.~Grinis reviewed the submission and reported, by e-mail, three concrete counter-examples that the existing suite failed to catch:

\begin{enumerate}
  \item \textbf{Option monad law violation.} \texttt{Option::pure(false)} returned \texttt{None} instead of \texttt{Some(false)}, because the implementation treated \texttt{false} as equivalent to \texttt{null}.
  \item \textbf{Monoid associativity violation.} For $a=0.004$, $b=0.004$, $c=0.005$ (euros), \texttt{TripMetrics} produced $(a\oplus b)\oplus c = 0.02$ but $a\oplus(b\oplus c)=0.01$.
  \item \textbf{Swallowed composition failure.} In \texttt{Itinerary::advance()}, a non-composable graph edge (target of the current path $\neq$ source of the next edge) was silently ignored instead of failing the itinerary computation.
\end{enumerate}

This report explains, for each issue, why the existing law was actually violated despite the suite reporting success, the minimal code change that restores the law, and the new regression check added so the suite can no longer miss it. The fixes are implemented directly in \texttt{src/CategoryTheory.php} and, where relevant, in the one call site in \texttt{index.php} that depended on the old (incorrect) behaviour.


\newpage
\section{Theoretical Background}\label{sec:background}
The three defects discussed in this report are violations of standard algebraic laws from category theory. This section fixes notation and restates the relevant definitions precisely enough to make the violations, and the proofs that the fixes restore the laws, unambiguous. The definitions follow Barr and Wells~\cite{BarrWells1990}; the monadic treatment of \texttt{Option}/\texttt{Either} follows Moggi~\cite{Moggi1991} and Wadler~\cite{Wadler1992}.

\subsection{Categories, Functors, and Natural Transformations}

A \emph{category} $\mathcal{C}$ consists of a collection of objects $\mathrm{Ob}(\mathcal{C})$, and for every pair of objects $u,v$ a set of morphisms $\mathrm{Hom}(u,v)$, together with an associative composition operator $\circ$ and, for every object $u$, an identity morphism $\mathrm{id}_u \in \mathrm{Hom}(u,u)$ satisfying
\[
  \mathrm{id}_v \circ f = f = f \circ \mathrm{id}_u \qquad \text{for all } f \in \mathrm{Hom}(u,v),
\]
\[
  (h \circ g) \circ f = h \circ (g \circ f) \qquad \text{whenever the composites are defined.}
\]
In \texttt{CategoryTheory.php}, the directed travel graph $G=(V,E)$ induces exactly such a category $F(G)$ (Section~\ref{sec:bug3} develops this in detail): $\mathrm{Ob}(F(G))=V$, and $\mathrm{Hom}(u,v)$ is the set of finite, directed walks from $u$ to $v$ along edges of $E$, with composition given by path concatenation and $\mathrm{id}_u$ the empty walk at $u$. This is the \emph{free category} generated by $G$: it has no morphisms beyond those forced to exist by the graph and the two axioms above.

A \emph{functor} $F : \mathcal{C} \to \mathcal{D}$ maps every object $u$ of $\mathcal{C}$ to an object $F(u)$ of $\mathcal{D}$ and every morphism $f : u \to v$ to a morphism $F(f) : F(u) \to F(v)$, subject to
\[
  F(\mathrm{id}_u) = \mathrm{id}_{F(u)}, \qquad F(g \circ f) = F(g) \circ F(f).
\]
\texttt{RouteList} implements this directly for the category of PHP values under function application: \texttt{RouteList::fmap(\$f)} maps every element with \texttt{\$f}, and \texttt{category\_laws\_test.php} checks both equations against the \texttt{discount}/\texttt{tax} transformations used throughout the booking pipeline.

A \emph{natural transformation} $\eta : F \Rightarrow G$ between two functors $F,G:\mathcal{C}\to\mathcal{D}$ is a family of morphisms $\eta_u : F(u) \to G(u)$, one per object $u$ of $\mathcal{C}$, such that for every $f : u \to v$ the naturality square commutes:
\[
  G(f) \circ \eta_u = \eta_v \circ F(f).
\]
\texttt{RouteList::headOption} is such a transformation from the \texttt{RouteList} functor to the \texttt{Option} functor: picking the head of a list and then mapping a function over the resulting \texttt{Option} must agree with mapping the function over the whole list first and then picking the head (both results are checked for equality in the suite, including the empty-list edge case where both sides must collapse to \texttt{None}).

\subsection{Monads, the Kleisli Category, and \texttt{Option}/\texttt{Either}}

A \emph{monad} on $\mathcal{C}$ is a functor $M$ together with two operations, $\mathrm{pure} : a \to M(a)$ and $\mathrm{bind} : M(a) \to (a \to M(b)) \to M(b)$ (written infix as $\mathbin{>\!\!>\!\!=}$), satisfying three laws for all $a$, $m : M(a)$ and Kleisli arrows $f,g$:
\[
  \text{(Left identity)}\quad \mathrm{pure}(a) \mathbin{>\!\!>\!\!=} f \;=\; f(a),
\]
\[
  \text{(Right identity)}\quad m \mathbin{>\!\!>\!\!=} \mathrm{pure} \;=\; m,
\]
\[
  \text{(Associativity)}\quad (m \mathbin{>\!\!>\!\!=} f) \mathbin{>\!\!>\!\!=} g \;=\; m \mathbin{>\!\!>\!\!=} \big(\lambda x.\, f(x) \mathbin{>\!\!>\!\!=} g\big).
\]
These three laws are exactly what make Kleisli arrows $a \to M(b)$ composable into a category (the \emph{Kleisli category} of $M$), with $\mathrm{pure}$ as identity and $\mathbin{>\!\!>\!\!=}$-chaining as composition. \texttt{Option} (the \texttt{Some}/\texttt{None} pair, PHP's analogue of \texttt{Maybe}) and \texttt{Either} (\texttt{Right}/\texttt{Left}, used for the short-circuiting validation pipeline in \texttt{index.php}) are both implemented as such monads. \textbf{The right-identity law is precisely the law violated by the first defect}: if $\mathrm{pure}$ does not act as a faithful identity on every value of the carrier type -- in particular on \texttt{bool} -- then $m \mathbin{>\!\!>\!\!=} \mathrm{pure}$ can produce a different \texttt{Option} than $m$ itself, and the Kleisli category stops being a category at all.

\subsection{Monoids and Strict Monoidal Functors}

A \emph{monoid} $(M, \oplus, e)$ is a set $M$ with an associative binary operation $\oplus$ and an identity element $e$:
\[
  (a \oplus b) \oplus c = a \oplus (b \oplus c), \qquad e \oplus a = a = a \oplus e.
\]
\texttt{TripMetrics} is declared a \texttt{Monoid} over $(\text{cost}, \text{duration}, \text{points})$ with \texttt{TripMetrics::empty()} as $e$ and \texttt{append()} as $\oplus$; it is used to fold the per-step costs of a booked itinerary into a single total. \textbf{The associativity law is precisely the law violated by the second defect}: Section~\ref{sec:bug2} exhibits concrete values for which $(a\oplus b)\oplus c \neq a\oplus(b\oplus c)$ under the original implementation.

A monoid homomorphism $M : (\mathcal{C}, \circ) \to (M,\oplus)$ from a category's morphisms into a monoid, satisfying $M(g \circ f) = M(f) \oplus M(g)$ and $M(\mathrm{id}) = e$, is called a \emph{strict monoidal functor}. \texttt{FreeCategoryPath::evaluateMetrics()} is exactly this: it evaluates a composed path in $F(G)$ by folding the \texttt{TripMetrics} monoid over the path's edges, and \texttt{category\_laws\_test.php} checks that evaluating a composed path agrees with evaluating and then combining its parts -- a check that is only meaningful if the underlying monoid actually associates, which is why the fix to \texttt{TripMetrics} in Section~\ref{sec:bug2} is a precondition for this check to mean anything at all.

\subsection{Error Propagation as a Categorical Obligation}

Short-circuiting a computation on failure is not merely a convenient control-flow idiom here: it is the \texttt{Either} monad's \texttt{bind} doing exactly what the monad laws require of it. If a step in a Kleisli chain produces \texttt{Left}, every subsequent \texttt{bind} in the chain must propagate that \texttt{Left} unchanged, by the definition of \texttt{Either::bind}. \textbf{The third defect is, in this sense, not a violation of one of the three named laws above, but a violation of the obligation those laws impose}: \texttt{FreeCategoryPath::compose} correctly computes a \texttt{Left} when two paths do not form a connected walk in $F(G)$, but \texttt{Itinerary::advance()} -- a plain method, not itself a Kleisli arrow -- read that \texttt{Either} and then discarded the failure instead of propagating it, breaking the chain of custody that the rest of the pipeline relies on. Section~\ref{sec:bug3} develops this precisely.


\newpage
\section{Analysis of the Three Reported Defects}\label{sec:motivation}

\subsection{Option Monad: \texttt{pure} Does Not Satisfy the Right-Identity Law}\label{sec:bug1}

The \texttt{Option} monad is expected to satisfy the monad laws for \emph{every} value its type parameter can hold, including \texttt{bool}. The original implementation of \texttt{Option::pure} was:

\begin{lstlisting}[caption={Original \texttt{Option::pure} (src/CategoryTheory.php)}]
public static function pure(mixed $value): self {
    return ($value === null || $value === false)
        ? new None()
        : new Some($value);
}
\end{lstlisting}

Taking $m = \mathtt{Some(false)}$ and $f = \texttt{Option::pure}$, the right-identity law $m \mathbin{>\!\!>\!\!=} \texttt{pure} \equiv m$ requires $m\mathbin{>\!\!>\!\!=}f$ to equal \texttt{Some(false)}. Instead it produced \texttt{None}, because \texttt{pure} conflated the domain value \texttt{false} with the \emph{absence} of a value — exactly the distinction the \texttt{Option} monad exists to preserve. The root cause is that \texttt{null} (absence) and \texttt{false} (a legitimate boolean payload) are not the same value, yet the guard treated them identically.

\textbf{Fix.} Only \texttt{null} denotes absence:
\begin{lstlisting}[caption={Fixed \texttt{Option::pure}}]
public static function pure(mixed $value): self {
    return ($value === null) ? new None() : new Some($value);
}
\end{lstlisting}

\textbf{Correctness argument.} For any $v \neq \mathrm{null}$, $\mathrm{pure}(v) = \mathrm{Some}(v)$ and $\mathrm{Some}(v)\mathbin{>\!\!>\!\!=}\mathrm{pure} = \mathrm{pure}(v) = \mathrm{Some}(v)$ by the definition of \texttt{Some::bind}, so right identity holds for every non-null value, including \texttt{false} -- the one case the original code broke. For $v=\mathrm{null}$, $\mathrm{pure}(\mathrm{null})=\mathrm{None}$ and $\mathrm{None}\mathbin{>\!\!>\!\!=}f=\mathrm{None}$ for any $f$ by \texttt{None::bind}, so the law holds there too. The fix restores the law for the entire domain, not just the reported counter-example.

One caller, \texttt{fetchLiveJsonOption()} in \texttt{index.php}, had silently relied on the old double meaning: \texttt{file\_get\_contents()} returns the sentinel \texttt{false} on failure, and the code built an \texttt{Option} from it via \texttt{Option::pure(\$raw)}. Since \texttt{pure} no longer treats \texttt{false} as absence, that call site was changed to test the sentinel explicitly instead of overloading a general-purpose constructor with a local, API-specific meaning:
\begin{lstlisting}[caption={index.php -- explicit sentinel check instead of overloaded pure()}]
$raw = @file_get_contents($url, false, $ctx);
if ($raw === false) return new None();
$decoded = json_decode($raw, true);
return $decoded === null ? new None() : new Some($decoded);
\end{lstlisting}
This call site, and every other place in the project that constructs or consumes an \texttt{Option}, is examined in full in the audit of Section~\ref{sec:audit}, and the fix is exercised against the two real external services it talks to in Section~\ref{sec:verification}.

\subsection{Monoid: \texttt{TripMetrics::append} Is Not Associative}\label{sec:bug2}

\texttt{TripMetrics} is declared a \texttt{Monoid}, which requires $(a\oplus b)\oplus c = a \oplus (b\oplus c)$ for all values, exactly. The original \texttt{append} rounded the running total to two decimals after \emph{every} addition:
\begin{lstlisting}[caption={Original \texttt{TripMetrics::append}}]
public function append(Monoid $other): self {
    return new self(
        round($this->costEur + $other->costEur, 2),
        ...
    );
}
\end{lstlisting}

For $a=0.004,\,b=0.004,\,c=0.005$ (Table~\ref{tabl:monoid}), rounding after each step amplifies the order in which values are combined into a visible difference: grouping $(a\oplus b)$ first rounds $0.008\to 0.01$, then $0.01+0.005\to 0.02$; grouping $(b\oplus c)$ first rounds $0.009\to 0.01$, then $a+0.01\to 0.01$. The two totals disagree, so the monoid is not actually associative — the repeated rounding is a lossy, order-dependent operation disguised as a pure sum.

\begin{table}[!ht]\centering
\caption{Associativity counter-example reported by Prof.~Grinis}
\label{tabl:monoid}
\begin{tabular}{|l|r|r|}\hline
Grouping & Original (buggy) & Fixed \\\hline
$(a\oplus b)\oplus c$ & 0.02 & 0.01 \\\hline
$a\oplus (b\oplus c)$ & 0.01 & 0.01 \\\hline
\end{tabular}
\end{table}

\textbf{Fix.} Round to whole cents exactly once, at construction, and keep that integer as the value actually added:
\begin{lstlisting}[caption={Fixed \texttt{TripMetrics} -- integer cents as the additive carrier}]
public readonly float $costEur;
public readonly int $costCents;

public function __construct(float $costEur, ...) {
    $this->costCents = (int) round($costEur * 100);
    $this->costEur = $this->costCents / 100;
}

public function append(Monoid $other): self {
    return new self(
        ($this->costCents + $other->costCents) / 100,
        ...
    );
}
\end{lstlisting}

\textbf{Correctness argument.} Let $c(x)$ denote the integer \texttt{costCents} stored by a \texttt{TripMetrics} value $x$, and let $x \oplus y$ denote \texttt{\$x->append(\$y)}. By construction, $c(x\oplus y) = c(x) + c(y)$ exactly, because \texttt{append} adds two already-rounded integers and the result's constructor round-trips an integer-valued cent amount back to the same integer (\texttt{round} of an exact or near-exact integer is a no-op up to floating-point noise far below one cent). Integer addition over $\mathbb{Z}$ is associative, so $c\big((a\oplus b)\oplus c\big) = c(a)+c(b)+c(c) = c\big(a\oplus(b\oplus c)\big)$ for \emph{all} $a,b,c$, not just the reported triple; since \texttt{costEur} is a deterministic function of \texttt{costCents} ($\mathrm{costEur}=c/100$), the two groupings also agree on \texttt{costEur}. No step in the computation rounds a value that has already absorbed a previous rounding, which is exactly what made the original version order-dependent.

\subsection{Free Category $F(G)$: Composition Failures Were Swallowed}\label{sec:bug3}

In the free category $F(G)$ induced by the travel graph, \texttt{FreeCategoryPath::compose} already returns an \texttt{Either}: \texttt{Right} on success, \texttt{Left} when the target vertex of the current path does not match the source vertex of the next edge (the paths do not form a connected walk). \texttt{Itinerary::advance()}, however, used that result incorrectly:
\begin{lstlisting}[caption={Original \texttt{Itinerary::advance} -- Left is discarded}]
$composed = $newPath->compose(FreeCategoryPath::fromEdge($graphEdge));
if ($composed->isRight()) {
    $newPath = $composed->getValue();
}
// execution continues either way
\end{lstlisting}
When composition fails, the \texttt{if} is simply not entered and \texttt{\$newPath} keeps its previous value; \texttt{advance()} still returns a brand-new, successful \texttt{Itinerary}. The categorical error is computed correctly and then thrown away, so an itinerary whose graph path is not actually a connected walk in $F(G)$ is reported as valid.

\textbf{Fix.} Propagate the failure instead of absorbing it:
\begin{lstlisting}[caption={Fixed \texttt{Itinerary::advance} -- Left fails the computation}]
$composed = $newPath->compose(FreeCategoryPath::fromEdge($graphEdge));
if (!$composed->isRight()) {
    $err = $composed->getValue();
    throw new RuntimeException("{$err['title']}: {$err['msg']}");
}
$newPath = $composed->getValue();
\end{lstlisting}

\textbf{Correctness argument.} After the fix, every execution path through the modified branch either (a) returns normally with \texttt{\$newPath} bound to the \texttt{Right} value of \texttt{compose} -- a path that is, by the definition of \texttt{compose}, a genuine morphism of $F(G)$ from the itinerary's start vertex to the new edge's target -- or (b) throws, and no \texttt{Itinerary} is constructed at all. There is no longer a third outcome where the method returns successfully with a path that \texttt{compose} itself rejected. A non-composable edge is an internal invariant violation of the step-building pipeline (every edge is constructed by the application itself from the previous step's state), not a user-facing validation outcome like a late flight or an over-budget trip, so failing fast with an exception is the appropriate signal, consistent with how every other \texttt{Left} produced inside this pipeline is already treated as a short-circuit.

\subsection{Regression Coverage}\label{sec:coverage}

Each fix above is paired with one new assertion in \texttt{tests/category\_laws\_test.php}, reproducing Prof.~Grinis' exact examples: \texttt{Some(false)} round-tripping through \texttt{bind(Option::pure)}; the $a=0.004,\,b=0.004,\,c=0.005$ monoid grouping; and a two-step \texttt{Itinerary::advance()} call with a deliberately mismatched edge, asserting that it throws. The full code of all three assertions is reproduced in Appendix~\ref{app:tests}, and running \texttt{php tests/category\_laws\_test.php} against the corrected code confirms all three now pass. Separately, the project's additional invariant suite, \texttt{LietuvaTravel\_regression\_checks.php}, exercises the \texttt{Option} monad, the \texttt{TripMetrics} monoid, and the free-category path/itinerary machinery under further edge-case inputs (nulls, sparse arrays, mismatched graph edges) and reports \textbf{10/10} checks passing.

\newpage
\section{Live Verification Against External Services}\label{sec:verification}
The three unit-level regression assertions added in Section~\ref{sec:coverage} verify the laws in isolation, with values constructed directly in PHP. They do not, by themselves, prove that the first fix behaves correctly at its real point of use: \texttt{fetchLiveJsonOption()} wraps \texttt{file\_get\_contents()} against two live third-party services (Open-Meteo for weather, Frankfurter/ECB for currency conversion), and the sentinel value that originally motivated the \texttt{Option::pure(false)} shortcut -- \texttt{file\_get\_contents()} returning \texttt{false} -- is produced by the PHP runtime itself on a real network failure, not by a value the test suite controls. To close that gap, the fixed application was started as an actual HTTP service and driven with real requests over the network.

\subsection{Method}

The built-in PHP development server was started against the corrected codebase, and the main page was requested twice: once with default parameters (exercising the live Open-Meteo weather call through \texttt{getLiveLithuanianWeather()}), and once with the currency explicitly set to \texttt{USD} (exercising the live Frankfurter call through \texttt{getLiveCurrencyRate()}, which is skipped entirely when the currency is \texttt{EUR}). Both calls route through the fixed \texttt{fetchLiveJsonOption()}:
\begin{lstlisting}[caption={Call site under live test (index.php)}]
function fetchLiveJsonOption(string $url): Option {
    $ctx = stream_context_create([...]);
    $raw = @file_get_contents($url, false, $ctx);
    if ($raw === false) return new None();
    $decoded = json_decode($raw, true);
    return $decoded === null ? new None() : new Some($decoded);
}
\end{lstlisting}

\begin{lstlisting}[language=bash,caption={Commands issued against the running server}]
php -S 127.0.0.1:8099 &
curl -s -o /dev/null -w "HTTP:%{http_code}\n" \
     http://127.0.0.1:8099/index.php

curl -s -X POST --data "currency=USD" \
     http://127.0.0.1:8099/index.php -o resp_usd.html
grep -oE "LIVE ECB|Fallback|LIVE Open-Meteo|Offline Cache" resp_usd.html
grep -iE "Fatal error|TypeError|Warning" resp_usd.html
\end{lstlisting}

\subsection{Result}

The default request returned \texttt{HTTP 200}. The \texttt{USD} request's response body contained \texttt{LIVE Open-Meteo} (twice -- once per weather read on the page) and, in the currency-exchange banner, a genuine live rate rather than the offline fallback:
\begin{lstlisting}
<strong>ECB Exchange Rate:</strong>
<span>1 EUR = 1.1224 USD</span>
\end{lstlisting}
No occurrence of \texttt{Fatal error}, \texttt{TypeError}, \texttt{Offline Cache}, or \texttt{Fallback} appeared in either response body or in the PHP development server's own log, which recorded clean \texttt{[200]} responses for both requests:
\begin{lstlisting}[language=bash]
[...] 127.0.0.1:58580 Accepted
[...] 127.0.0.1:58580 [200]: POST /index.php
[...] 127.0.0.1:58580 Closing
\end{lstlisting}

\subsection{Why This Matters}

Before the fix, \texttt{Option::pure(\$raw)} treated a \texttt{false} \texttt{\$raw} (a failed fetch) and a string \texttt{\$raw} (a successful fetch) identically whenever the application code called \texttt{pure} directly on the raw result -- which it previously did. Removing the \texttt{false}-as-absence special case from \texttt{Option::pure} without touching this call site would have made a \emph{successful} fetch behave correctly (since \texttt{Option::pure(\$raw)} would now wrap any non-null string as before) but would have broken a \emph{failed} fetch in a new way: \texttt{Option::pure(false)} would now produce \texttt{Some(false)} instead of \texttt{None}, and a subsequent \texttt{fmap} expecting a decoded payload would be invoked against that \texttt{false} value instead of short-circuiting -- turning a graceful offline fallback into a runtime error on every network hiccup. This is exactly the kind of regression that a purely static fix, verified only against hand-picked values, can introduce without anyone noticing until the live service degrades. The explicit sentinel check shown above (\texttt{if (\$raw === false) return new None();}, followed by an explicit \texttt{null}-check on the decoded JSON so malformed responses also degrade to \texttt{None} rather than \texttt{Some(null)}) was therefore written and then exercised against the two real external services before this report was finalised, rather than left as an untested consequence of the monad-law fix.


\newpage
\section{Audit of All Call Sites Depending on \texttt{Option::pure}}\label{sec:audit}
Changing the semantics of a foundational primitive such as \texttt{Option::pure} is, by construction, a project-wide change: every caller that happened to rely on the old (incorrect) double meaning of \texttt{false} is now exposed to different behaviour. Fixing the one call site found by inspection (\texttt{fetchLiveJsonOption}, Section~\ref{sec:verification}) is not sufficient evidence that no other call site was silently depending on the bug; it is evidence for that one call site only. An explicit audit of the entire codebase was therefore run before considering the fix complete.

\begin{lstlisting}[language=bash,caption={Audit command: every construction and consumption site of Option}]
grep -rn "Option::pure\|->fmap(\|->bind(\|->getOrElse(\|new Some(\|new None(" \
     index.php src/
\end{lstlisting}

The search returns every place in the project that either constructs an \texttt{Option} (via \texttt{Option::pure}, \texttt{new Some}, or \texttt{new None}) or consumes one (via \texttt{fmap}, \texttt{bind}, or \texttt{getOrElse}). Two structurally different kinds of result came back:

\begin{enumerate}
  \item \textbf{Direct, explicit construction} -- \texttt{new Some(\dots)} and \texttt{new None()} appear directly in \texttt{RouteList::headOption} (natural transformation on an empty/non-empty list), in \texttt{Flight}/\texttt{Hotel}'s fallback offers, and in the now-fixed \texttt{fetchLiveJsonOption}. None of these ever pass through \texttt{Option::pure}, so none of them were affected by its old \texttt{false}-collapsing behaviour, and none of them are affected by removing it.
  \item \textbf{Indirect construction via \texttt{Option::pure}} -- exactly \emph{two} call sites in the entire codebase invoke \texttt{Option::pure} itself: \texttt{Some::fmap()} (line~39 of \texttt{CategoryTheory.php}, which builds its result as \texttt{Option::pure(\$f(\$this->value))}) and the original \texttt{fetchLiveJsonOption} (already fixed). No third call site exists.
\end{enumerate}

\texttt{Some::fmap} deserves a second look, because it was not the call site Prof.~Grinis reported, yet it shares the same root cause: before the fix, mapping a function that returns \texttt{false} over a \texttt{Some} value -- e.g.\ \texttt{Some(5)->fmap(fn(\$x) => \$x > 10)} -- silently produced \texttt{None} instead of \texttt{Some(false)}, violating the functor identity law in spirit (a successful computation whose result happens to be \texttt{false} should not disappear). The single-line fix to \texttt{Option::pure} repairs \texttt{Some::fmap} automatically, since it is defined in terms of \texttt{pure}; no separate change to \texttt{fmap} was necessary, and the audit confirms there is no third, undiscovered call site that would need one.

\textbf{Conclusion of the audit.} The fix to \texttt{Option::pure} has exactly two points of impact in the whole project -- \texttt{Some::fmap} (now correct, no change needed) and \texttt{fetchLiveJsonOption} (fixed and verified live in Section~\ref{sec:verification}). No other code in \texttt{index.php} or \texttt{src/CategoryTheory.php} constructs an \texttt{Option} through \texttt{pure}, so no other behaviour in the application changes as a result of this fix.


 %Conclusions section
\sectionWithoutNumber{\keyWordConclusions}{conclu}
All three defects reported by Prof.~I.~Grinis were root-caused and fixed with minimal, localised changes to \texttt{src/CategoryTheory.php} (and one dependent call site in \texttt{index.php}):

\begin{itemize}
  \item \texttt{Option::pure} now treats only \texttt{null} as absence of value, so \texttt{Some(false)} round-trips correctly through \texttt{bind(Option::pure)}, restoring the monad's right-identity law. The single call site that relied on the old behaviour (\texttt{fetchLiveJsonOption}, which used the \texttt{false} sentinel returned by \texttt{file\_get\_contents} on failure) was updated to test for that sentinel explicitly instead of overloading \texttt{Option::pure} with two meanings.
  \item \texttt{TripMetrics} now rounds to whole cents exactly once, at construction time, and stores that integer value (\texttt{costCents}). \texttt{append()} adds integers, which is exactly associative, instead of adding floats and re-rounding to two decimals after every step, which was not. The public float interface (\texttt{costEur}) is unchanged, so no other call site required modification.
  \item \texttt{Itinerary::advance()} now throws a \texttt{RuntimeException} carrying the free category's own \texttt{Left} error message when a graph edge does not compose with the current induced path, instead of silently keeping the previous (shorter, incorrect) path and returning a successful itinerary.
\end{itemize}

Each fix is accompanied by a new assertion in \texttt{tests/category\_laws\_test.php}, reproducing the exact counter-example supplied by Prof.~Grinis (including the $a=0.004,\,b=0.004,\,c=0.005$ monoid case). The three new assertions fail against the old code and pass against the fixed code, confirming both the regressions and the fixes. Separately, the project's additional invariant suite, \texttt{LietuvaTravel\_regression\_checks.php}, reports \textbf{10/10} checks passing against the corrected code. Table~\ref{tabl:summary} summarises the three defects end to end.

\begin{table}[!ht]\centering
\caption{Summary of the three defects, from report to fix to verification}
\label{tabl:summary}
\small
\begin{tabular}{|p{2.6cm}|p{3.6cm}|p{4.3cm}|p{3.3cm}|}\hline
\textbf{Law violated} & \textbf{Root cause} & \textbf{Fix} & \textbf{Verified by} \\\hline
Monad right identity ($m\mathbin{>\!\!>\!\!=}\mathrm{pure}\equiv m$) & \texttt{false} treated as absence in \texttt{Option::pure} & Only \texttt{null} denotes absence; sentinel checked explicitly at the one real caller & Unit assertion + live HTTP test (Sec.~\ref{sec:verification}) \\\hline
Monoid associativity ($(a\oplus b)\oplus c = a\oplus(b\oplus c)$) & Re-rounding to 2 decimals after every \texttt{append} & Accumulate in integer cents, round once at construction & Unit assertion with Prof.~Grinis' exact $a,b,c$ \\\hline
Error-propagation obligation of \texttt{Either::bind} & \texttt{Left} from \texttt{compose} read but not propagated in \texttt{advance()} & Throw on \texttt{Left}, carrying the categorical error message & Unit assertion asserting the throw occurs \\\hline
\end{tabular}
\end{table}

This progression -- a law-level description of what went wrong, a root-cause diagnosis in the code, a minimal fix, and an independent check that the fix actually holds -- was applied uniformly to all three defects, and is the same discipline the project's own law suite was supposed to enforce in the first place. The suite's earlier ``15/15'' badge was not dishonest, only incomplete: it verified every law it happened to test, but never generated the adversarial inputs (a falsy payload, a sub-cent amount, a deliberately broken graph) that expose these particular violations. The three new assertions close exactly that gap.


\sectionWithoutNumber{Future Work}{future}{}
Beyond the three fixes above, two concrete follow-ups would further harden this codebase's categorical guarantees. First, the remaining ad hoc \texttt{InvalidArgumentException} checks inside \texttt{Some::bind} and \texttt{Right::bind} (``must return an Option/Either instance'') are runtime guards against a mistake the type system could catch instead; PHP's lack of higher-kinded types makes a fully static alternative awkward, but a static-analysis baseline (e.g.\ Psalm or PHPStan with a custom plugin recognising the Kleisli-arrow shape) could flag a non-conforming callback at the call site rather than at the moment it executes. Second, the law suite itself tests fixed, hand-picked values; a property-based layer (generating random floats, booleans, and graph shapes, and checking each law automatically over hundreds of generated cases per run) would very likely have caught \texttt{Option::pure(false)} and the sub-cent monoid rounding without waiting for a human reviewer to supply the exact counter-example, and is the most direct way to prevent a fourth such defect from reaching review.

\sectionWithoutNumber{\keyWordAIdeclaration}{aideclare}
Claude Code (Anthropic), an AI coding assistant, was used under the author's direction to: (i) locate the exact lines responsible for each of the three defects reported by Prof.~Grinis; (ii) implement the corresponding fixes in \texttt{src/CategoryTheory.php} and \texttt{index.php}; (iii) add the three regression assertions to \texttt{tests/category\_laws\_test.php}; (iv) run the test suite (\texttt{php tests/category\_laws\_test.php}) to confirm all laws pass, and the additional invariant suite (\texttt{php LietuvaTravel\_regression\_checks.php}) to confirm 10/10 checks pass; and (v) draft this LaTeX report from the official department template. All changes were reviewed by the author before submission. No AI-generated content was included without verification against the actual code and test output.


 %file literatureSources.bib
\referenceSources{literatureSources}

\newpage
\begin{appendices}
\section{Corrected Core Classes (\texttt{src/CategoryTheory.php})}\label{app:src}

% Listings below carry the file's real (3-digit) line numbers via firstnumber,
% so the numbering column needs to be widened or the digits get clipped by the frame.
\lstset{xleftmargin=32pt, framexleftmargin=24pt, numbersep=5pt}

For reference, this appendix reproduces -- verbatim, with line numbers taken from the file as submitted -- the four classes directly touched by the three fixes: the \texttt{Option} monad (Listing~\ref{lst:app-option}), the \texttt{TripMetrics} monoid (Listing~\ref{lst:app-monoid}), the free category path type \texttt{FreeCategoryPath} (Listing~\ref{lst:app-path}), and \texttt{Itinerary::advance()} (Listing~\ref{lst:app-advance}).

\begin{lstlisting}[caption={\texttt{Option}, \texttt{Some}, \texttt{None} -- fixed \texttt{pure}},label={lst:app-option},firstnumber=27]
abstract class Option implements Monad {
    public static function pure(mixed $value): self {
        return ($value === null) ? new None() : new Some($value);
    }
    abstract public function isSome(): bool;
    abstract public function getOrElse(mixed $default): mixed;
    abstract public function toEither(array $errorData): Either;
}

final class Some extends Option {
    public function __construct(private readonly mixed $value) {}
    public function fmap(callable $f): Option { return Option::pure($f($this->value)); }
    public function bind(callable $f): Option {
        $res = $f($this->value);
        if (!$res instanceof Option) {
            throw new InvalidArgumentException("Option::bind must return an Option instance.");
        }
        return $res;
    }
    public function isSome(): bool { return true; }
    public function getOrElse(mixed $default): mixed { return $this->value; }
    public function toEither(array $errorData): Either { return new Right($this->value); }
}

final class None extends Option {
    public function fmap(callable $f): Option { return $this; }
    public function bind(callable $f): Option { return $this; }
    public function isSome(): bool { return false; }
    public function getOrElse(mixed $default): mixed { return $default; }
    public function toEither(array $errorData): Either { return new Left($errorData); }
}
\end{lstlisting}

\begin{lstlisting}[caption={\texttt{TripMetrics} -- fixed, cent-exact \texttt{append}},label={lst:app-monoid},firstnumber=122]
final class TripMetrics implements Monoid {
    public readonly float $costEur;
    public readonly int $costCents;

    public function __construct(
        float $costEur,
        public readonly int $durationHours,
        public readonly int $points,
        public readonly array $items = []
    ) {
        $this->costCents = (int) round($costEur * 100);
        $this->costEur = $this->costCents / 100;
    }

    public static function empty(): self { return new self(0.0, 0, 0, []); }

    public function append(Monoid $other): self {
        if (!$other instanceof self) {
            throw new InvalidArgumentException("Can only append TripMetrics.");
        }
        return new self(
            ($this->costCents + $other->costCents) / 100,
            $this->durationHours + $other->durationHours,
            $this->points + $other->points,
            array_merge($this->items, $other->items)
        );
    }

    public static function fold(array $list): self {
        return array_reduce($list, fn(self $acc, self $item) => $acc->append($item), self::empty());
    }
}
\end{lstlisting}

\begin{lstlisting}[caption={\texttt{FreeCategoryPath} -- composition in $F(G)$, unchanged},label={lst:app-path},firstnumber=183]
final class FreeCategoryPath {
    public function __construct(
        public readonly GraphNode $source,
        public readonly GraphNode $target,
        public readonly array $edges = []
    ) {}

    public static function identity(GraphNode $node): self {
        return new self($node, $node, []);
    }

    public static function fromEdge(GraphEdge $edge): self {
        return new self($edge->source, $edge->target, [$edge]);
    }

    public function compose(self $next): Either {
        if ($this->target->id !== $next->source->id) {
            return new Left([
                'step'  => 'Free Category Inductor F(G)',
                'title' => 'Non-Composable Graph Morphisms',
                'msg'   => "Target vertex '{$this->target->id}' does not match source vertex '{$next->source->id}'.",
                'fix'   => 'Ensure graph edges form a connected walk u -> v -> w.'
            ]);
        }
        return new Right(new self(
            $this->source,
            $next->target,
            array_merge($this->edges, $next->edges)
        ));
    }

    public function evaluateMetrics(): TripMetrics {
        return TripMetrics::fold(array_map(fn(GraphEdge $e) => $e->weight, $this->edges));
    }
}
\end{lstlisting}

\begin{lstlisting}[caption={\texttt{Itinerary::advance()} -- fixed, propagates composition failure},label={lst:app-advance},firstnumber=348]
public function advance(
    ?Flight $flight = null,
    ?Hotel $hotel = null,
    ?array $activities = null,
    ?TripMetrics $stepMetrics = null,
    string $stepType = '',
    string $traceMsg = '',
    ?GraphEdge $graphEdge = null
): self {
    $newMetrics = $stepMetrics ? $this->metrics->append($stepMetrics) : $this->metrics;
    $newTrace = $this->trace;
    if ($traceMsg !== '') $newTrace[] = $traceMsg;

    $newPath = $this->inducedPath;
    if ($graphEdge !== null) {
        if (empty($newPath->edges)) {
            $newPath = FreeCategoryPath::fromEdge($graphEdge);
        } else {
            $composed = $newPath->compose(FreeCategoryPath::fromEdge($graphEdge));
            if (!$composed->isRight()) {
                $err = $composed->getValue();
                throw new RuntimeException("{$err['title']}: {$err['msg']}");
            }
            $newPath = $composed->getValue();
        }
    }

    return new self(/* ... unchanged field threading ... */);
}
\end{lstlisting}

\section{New Regression Assertions (\texttt{tests/category\_laws\_test.php})}\label{app:tests}

Listing~\ref{lst:app-test1} through Listing~\ref{lst:app-test3} are the three assertions added to the suite, one per defect, each reproducing Prof.~Grinis' exact counter-example rather than a weaker, hand-picked one.

\begin{lstlisting}[caption={Regression for the \texttt{Option::pure(false)} defect},label={lst:app-test1},firstnumber=99]
assertLaw(
    '3. Option Monad Laws (Maybe)',
    'Pure Preserves Falsy Values (Right Identity, false case)',
    'm >>= Option::pure  ===  m  for m = Some(false)',
    (new Some(false))->bind($optPure)->isSome()
        && (new Some(false))->bind($optPure)->getOrElse('MISSING') === false,
    'Some(false)->bind(Option::pure) must stay Some(false), not collapse to None (false !== absence)'
);
\end{lstlisting}

\begin{lstlisting}[caption={Regression for the sub-cent monoid associativity defect},label={lst:app-test2},firstnumber=195]
$cent1 = new TripMetrics(0.004, 0, 0);
$cent2 = new TripMetrics(0.004, 0, 0);
$cent3 = new TripMetrics(0.005, 0, 0);
$centGroup1 = $cent1->append($cent2)->append($cent3);
$centGroup2 = $cent1->append($cent2->append($cent3));
assertLaw(
    '5. Monoid Laws (TripMetrics)',
    'Associativity Law (Sub-Cent Rounding Case)',
    '(a (+) b) (+) c  =  a (+) (b (+) c)  for a=0.004, b=0.004, c=0.005',
    $centGroup1->costEur === $centGroup2->costEur,
    "(EUR0.004 (+) EUR0.004) (+) EUR0.005 = EUR{$centGroup1->costEur} === EUR0.004 (+) (EUR0.004 (+) EUR0.005) = EUR{$centGroup2->costEur}"
);
\end{lstlisting}

\begin{lstlisting}[caption={Regression for the swallowed graph-composition-failure defect},label={lst:app-test3},firstnumber=263]
$it0 = new Itinerary('Vilnius', 'VNO', 2, 1, 500.0, 'Sunny');
$it1 = $it0->advance(graphEdge: new GraphEdge(
    'e0', new GraphNode('VNO', 'Origin Hub', 'origin'), $nVNO, 'Arrive', TripMetrics::empty()
));
$advanceThrew = false;
try {
    $it1->advance(graphEdge: new GraphEdge('bad', $nHTL, $nACT, 'Mismatch', TripMetrics::empty()));
} catch (RuntimeException $e) {
    $advanceThrew = true;
}
assertLaw(
    '7. Free Category Inductor F(G) (Itinerary Fail-Fast)',
    'Composition Failure Propagation',
    'compose(p1, p2) = Left  =>  advance() throws, never silently ignores it',
    $advanceThrew,
    'Itinerary::advance() throws RuntimeException when the induced path target does not match the next edge source.'
);
\end{lstlisting}

\end{appendices}

\end{document}
